\par\medskip
\noindent\textbf{Evidencias de la Fase 1: Análisis, prototipado y modelado de datos (mayo – junio de 2026)}
\par\nobreak
A partir de las entrevistas realizadas al Gerente General y al Contador, 
se evaluaron soluciones comerciales del mercado peruano (Buk Perú, Atenea Perú y Rapifac Nóminas) 
tomando a Rapifac como referencia operativa principal por la disponibilidad de flujos demostrativos completos. 
Este análisis permitió delimitar los 15 requerimientos funcionales (Anexo~2) y 9 no funcionales (Anexo~3) del MVP para los perfiles de RR.HH. y Contador, 
postergando el portal de autoservicio del trabajador como línea de evolución futura (Tabla~\ref{tab:comparativa_referentes}).

\par
\noindent
\begin{minipage}{\linewidth}

    \tituloTablaAPA
        {tab:comparativa_referentes}
        {Análisis comparativo de sistemas de referencia y aportes del MVP propuesto}

    \begingroup
    \renewcommand{\baselinestretch}{1}
    \footnotesize
    \setlength{\tabcolsep}{4pt}
    \renewcommand{\arraystretch}{1.35}

    \begin{tabularx}{\linewidth}{
        @{}
        >{\raggedright\arraybackslash}p{0.20\linewidth}
        >{\raggedright\arraybackslash}p{0.26\linewidth}
        >{\raggedright\arraybackslash}p{0.26\linewidth}
        >{\raggedright\arraybackslash}X
        @{}
    }
        \toprule
        \textbf{Sistema evaluado} &
        \textbf{Fortalezas observadas} &
        \textbf{Limitaciones para el caso} &
        \textbf{Aporte en el MVP de J\&P} \\
        \midrule

        Buk Perú &
        Suite integral corporativa de capital humano y nóminas. &
        Tarificación por colaborador activo y costo elevado para pymes. &
        Modelo por tramos con arquitectura modular escalable. \\

        Atenea Perú / Planiweb &
        Paquetes planos accesibles orientados a contabilidad pyme. &
        Menor flexibilidad en cruce de asistencia y flujos de aprobación. &
        Cruce automático de jornadas, marcaciones y solicitudes con sustento. \\

        Rapifac Nóminas &
        Cobertura normativa local (PLAME, AFPnet, CTS, gratificaciones, 5ta). &
        Alta carga manual en boletas individuales y formularios extensos. &
        Wizard de registro con RENIEC, carga masiva en dos pasos y envíos con n8n. \\

        \bottomrule
    \end{tabularx}
    \endgroup

    \notaTabla{Síntesis del levantamiento comparativo ejecutado en la Fase~1 de análisis. Elaboración propia.}

\end{minipage}
\finBloqueAPA

Como parte del diseño temprano de interfaz (actividad registrada en el cronograma de la Fase~1), 
se elaboraron prototipos exploratorios en Google Stitch para validar la disposición de indicadores, 
alertas contractuales y formularios de ingreso de personal antes de iniciar la codificación 
(Figuras~\ref{fig:stitch_dashboard_rrhh} a~\ref{fig:stitch_wizard_empleado}).

\figuraEnTexto
    {fig:stitch_dashboard_rrhh}
    {Prototipo inicial del panel de gestión de Recursos Humanos en Google Stitch}
    {images/evidencias/fase1/dashboard-rrhh.png}
    {Prototipo exploratorio generado con Google Stitch mediante el \textit{prompt}: 
    ``Panel administrativo de RR.HH. para J\&P Planillas usando componentes Shadcn y color primario \#150AB4, con tarjetas de empleados totales, solicitudes activas, documentos pendientes, alertas tempranas de contratos, gráfico por área y próximos cumpleaños''.}

\figuraEnTexto
    {fig:stitch_dashboard_planillas}
    {Prototipo inicial del panel de control de nóminas en Google Stitch}
    {images/evidencias/fase1/dashboard-planillas.png}
    {Prototipo exploratorio con datos referenciales generado con Google Stitch mediante el \textit{prompt}: ``Dashboard de planillas con estilo moderno en \#150AB4, cuatro tarjetas KPI superiores (total nómina, empleados, pendientes, retenciones), gráfico de cumplimiento de pago circular, alertas activas de PLAME/DNI/horas extras y barra de progreso mensual''.}

\figuraEnTexto
    {fig:stitch_wizard_empleado}
    {Prototipo del formulario en pasos para registro de colaboradores con consulta RENIEC}
    {images/evidencias/fase1/wizard-1.png}
    {Prototipo exploratorio generado con Google Stitch mediante el \textit{prompt}: ``Modal tipo asistente (\textit{wizard}) de tres pasos para registro de trabajador (Datos Personales, Información Laboral, Pago y Pensiones) con búsqueda automática de DNI en RENIEC''.}

Asimismo, el modelado relacional de la base de datos en PostgreSQL 
se estructuró por dominios funcionales para garantizar la integridad referencial, 
la inmutabilidad de los cierres contables y la trazabilidad en auditoría. Las Figuras~\ref{fig:erd_organizacion_legajos} 
a~\ref{fig:erd_seguridad} presentan las 23 tablas operativas del MVP agrupadas en cuatro diagramas modulares 
(conservando adicionalmente 7 tablas previstas para evolución futura de vínculos históricos, vacaciones y liquidaciones).

\figuraEnTexto
    {fig:erd_organizacion_legajos}
    {Modelo relacional del módulo de organización, legajos y contratos}
    {images/evidencias/fase1/ERD_1_organizacion_legajos.png}
    {Entidades que sostienen la estructura organizacional y el legajo del colaborador. PK = clave primaria; FK = clave foránea; UQ = valor único. 
    \texttt{solicitud.responsable\_id} referencia a \texttt{usuarios} (Figura~\ref{fig:erd_seguridad}). Elaboración propia.}

\figuraEnTexto
    {fig:erd_asistencia}
    {Modelo relacional del módulo de asistencia e incidencias}
    {images/evidencias/fase1/ERD_2_asistencia.png}
    {Registro de marcaciones, vigencias de jornada, decisiones operativas, calendario laboral y consolidado mensual. 
    \texttt{empleados} y \texttt{jornada} actúan como tablas puente; \texttt{usuario\_id} y \texttt{aprobado\_por} referencian a \texttt{usuarios} (Figura~\ref{fig:erd_seguridad}). Elaboración propia.}

\figuraEnTexto
    {fig:erd_nomina_parametros}
    {Modelo relacional del módulo de nómina, parámetros y datos financieros}
    {images/evidencias/fase1/ERD_3_nomina_parametros.png}
    {Datos bancarios y previsionales, catálogos de bancos y AFP, parámetros legales, novedades, adelantos y resultados de planilla. 
    \texttt{parametro\_legal} no posee clave foránea al ser consultada por código y vigencia desde el motor de cálculo. Elaboración propia.}

\figuraEnTexto
    {fig:erd_seguridad}
    {Modelo relacional del módulo de seguridad y operación}
    {images/evidencias/fase1/ERD_4_seguridad.png}
    {Cuentas de usuario vinculadas al legajo, tokens de seguridad, bitácora de auditoría, tareas operativas y trabajos de carga masiva. 
    \texttt{asignado\_a} y \texttt{asignado\_por} referencian a \texttt{usuarios}. Elaboración propia.}

\par\medskip
\noindent\textbf{Evidencias de la Fase 2: Construcción y validación del \textit{backend} (julio – octubre de 2026)}\par
\nobreak
Durante la Fase~2 se implementó el núcleo transaccional en NestJS~11, Fastify~5 y Prisma ORM sobre PostgreSQL, junto con los hilos paralelos de cálculo (\textit{Worker Threads}) y los disparadores SQL de inmutabilidad y auditoría. La evidencia de esta etapa se sustenta en cuatro artefactos técnicos verificables del repositorio \texttt{bella-jyp-backend}:

\begin{itemize}[leftmargin=*]
    \item \textbf{Trazabilidad de código y migraciones:} Historial de confirmaciones en GitHub entre julio y octubre de 2026 y registro de migraciones en Prisma (entre ellas \texttt{20261001...\_refactor\_schema\_afp\_planillas\_estados} y \texttt{20261015...\_asistencia\_payroll\_estable}), acreditando la evolución incremental de los cinco módulos de dominio (\texttt{core}, \texttt{RRHH}, \texttt{asistencia}, \texttt{nomina} e \texttt{integraciones}) (Figura~\ref{fig:evidencia_git_backend}).
    \item \textbf{Contratos de interfaz y pruebas de \textit{endpoints}:} Especificación formal \texttt{openapi.yml} (versión~1.2.0) publicada en Swagger UI y colecciones de Postman (\textit{Bella J\&P — Avance backend} y \textit{Planilla}), validando códigos de respuesta, esquemas Zod y manejo de errores bajo RFC~7807 (Figuras~\ref{fig:evidencia_swagger} y~\ref{fig:evidencia_postman}).
    \item \textbf{Pruebas automatizadas y análisis estático:} Ejecución de 180 archivos de prueba unitaria y de integración (\texttt{.spec.ts} y \texttt{.e2e-spec.ts}) en Jest sobre las reglas laborales y disparadores SQL, complementada con el análisis de calidad en SonarQube (\texttt{bella-planillas-backend}) conforme a la definición de Terminado (Figura~\ref{fig:evidencia_jest_sonar}).
\end{itemize}

%\figuraEnTexto
%    {fig:evidencia_git_backend}
%    {Historial de actividad y migraciones del repositorio \textit{backend}}
%    {images/evidencias/fase2/github-commits-backend.png}
%    {Registro de confirmaciones en el repositorio \texttt{bella-jyp-backend} correspondientes al desarrollo del núcleo transaccional (julio – octubre de 2026). Elaboración propia.}

%\figuraEnTexto
%    {fig:evidencia_swagger}
%    {Documentación interactiva de contratos API en Swagger UI (OpenAPI 3.0.3)}
%    {images/evidencias/fase2/swagger-openapi.png}
%    {Endpoints agrupados por dominio funcional a partir de la especificación \texttt{openapi.yml} v1.2.0. Elaboración propia.}

%\figuraEnTexto
%    {fig:evidencia_postman}
%    {Validación funcional de \textit{endpoints} de nómina en Postman}
%    {images/evidencias/fase2/postman-planilla.png}
%    {Ejecución de prueba de integración sobre la colección del \textit{backend} verificando respuesta JSON estructurada y códigos de estado HTTP. Elaboración propia.}

%\figuraEnTexto
%    {fig:evidencia_jest_sonar}
%    {Resultados de pruebas automatizadas en Jest y análisis en SonarQube}
%    {images/evidencias/fase2/jest-sonarqube.png}
%    {Verificación de cobertura de pruebas unitarias e integración y ausencia de defectos bloqueantes en el proyecto \texttt{bella-planillas-backend}. Elaboración propia.}


\par\medskip
\noindent\textbf{Evidencias de las Fases 3, 4 y 5: Integración \textit{frontend}, automatización con n8n y despliegue (noviembre – diciembre de 2026)}\par
\nobreak